% !TEX spellcheck = en_US



% ~~~~~~~~~~~~~~~~~~~~~~~~~~~~~~~~~~~~~~~ V
% (NC) 2026 Haluk Bingol
% github.com/halukbingol/LaTeX-Templates
%
% Licensees may copy, distribute, display, and perform the work and make derivative works 
% and remixes based on it only for non-commercial purposes. 
% ~~~~~~~~~~~~~~~~~~~~~~~~~~~~~~~~~~~~~~~ A




% =====================================================================
%  The Eclipse IDE -- Sophomore Level: Projects, Tests, Debugging and Refactoring
%  ---------------------------------------------------------------------
%  Built on hblib/Template/hblibTemplate.tex with the libraries in ../../hblib/.
%  Programs and their real outputs are in codeoutput/:
%     codeoutput/<Name>.java  a program (classes without main() are only compiled)
%     codeoutput/<Name>.sh    a terminal session, run in a copy of codeoutput/<Name>/
%     codeoutput/<Name>.txt   what it printed, made by:  bash run_examples.sh
% =====================================================================

\documentclass[aspectratio=169,10pt,compress]{beamer}

% ~~~~~~~~~~~~~~~~~~~~~~~~~~~~~~~~~~~~~~~~~~~~~~~~~~~~~~~~~~~~~~~~~ V hblib
	\usepackage{../../hblib/hblibPresentation} % lib presentation: theme, i/N, \hSection, algorithm2e, ...
	\usepackage{../../hblib/hblibCodeoutput} % lib code + output: \lstcode, \lstterminal, ...
	\usepackage{../../hblib/hblibDatastructures} % lib data structures: \hString, \hStack, \hQueue
% ~~~~~~~~~~~~~~~~~~~~~~~~~~~~~~~~~~~~~~~~~~~~~~~~~~~~~~~~~~~~~~~~~ A hblib


% theme, i/N, fonts (T1, lmodern), booktabs, xcolor, array (L{width} columns), listings, tikz,
% algorithm2e (\small captions, \Fn) and the color gitred come from the hblib libraries




% +++++++++++++++++++++++++++++++++++++++ V specific
\DontPrintSemicolon % in algorithms
\newcommand{\win}{\textbf{Windows}}
\newcommand{\mac}{\textbf{macOS}}
% +++++++++++++++++++++++++++++++++++++++ A specific

% ---------------------------------------------------------------------
\title{The Eclipse IDE}
\subtitle{Sophomore Level: Projects, Tests, Debugging and Refactoring}
\author[Bingol]{Haluk O. Bingol}

\institute[FCIS]{
    Faculty of Computer and Information Sciences\\
    Yeditepe University
}
\date{
	{\tiny \hbTimeStamp}
}

\begin{document}




% ~~~~~~~~~~~~~~~~~~~~~~~~~~~~~~~~~~~~~~ V frm
\begin{frame}[t,fragile]
  \titlepage
\end{frame}
% ~~~~~~~~~~~~~~~~~~~~~~~~~~~~~~~~~~~~~~ A frm




% ~~~~~~~~~~~~~~~~~~~~~~~~~~~~~~~~~~~~~~ V frm
\begin{frame}[t,fragile] \frametitle{Outline}
  \begin{columns}[T]
    \begin{column}{0.48\textwidth}
      \tableofcontents[sections={1-4}, hideallsubsections]
    \end{column}
    \begin{column}{0.48\textwidth}
      \tableofcontents[sections={5-8}, hideallsubsections]
    \end{column}
  \end{columns}
\end{frame}
% ~~~~~~~~~~~~~~~~~~~~~~~~~~~~~~~~~~~~~~ A frm

% ~~~~~~~~~~~~~~~~~~~~~~~~~~~~~~~~~~~~~~ V
\hSection[Projects]{Projects with Packages}
% ~~~~~~~~~~~~~~~~~~~~~~~~~~~~~~~~~~~~~~ A




% ~~~~~~~~~~~~~~~~~~~~~~~~~~~~~~~~~~~~~~ V frm
\begin{frame}[t,fragile] \frametitle{A project with packages}
  \begin{columns}[T]
    \begin{column}{0.5\textwidth}
      \begin{itemize}\small
        \item \menu{File > New > Package}: \texttt{tr.edu.yeditepe.grades}. A package is a \textbf{folder}:
              \texttt{src/tr/edu/yeditepe/grades/}.
        \item Package names: reversed domain, lower case. They prevent name clashes between libraries.
        \item Every file starts with its \texttt{package} line; Eclipse writes it and keeps it right when you
              move a class (\menu{Refactor > Move}).
        \item Package Explorer shows packages flat; \menu{Window > Show View > Navigator} shows the real folders.
      \end{itemize}
    \end{column}
    \begin{column}{0.46\textwidth}
      \lstcode[basicstyle=\ttfamily\scriptsize]{codeoutput/p/Student.java}{Java}
      \lstcode[basicstyle=\ttfamily\tiny, firstline=13, firstnumber=13]{codeoutput/p/GradeBook.java}{Java}
    \end{column}
  \end{columns}
\end{frame}
% ~~~~~~~~~~~~~~~~~~~~~~~~~~~~~~~~~~~~~~ A frm




% ~~~~~~~~~~~~~~~~~~~~~~~~~~~~~~~~~~~~~~ V frm
\begin{frame}[t,fragile] \frametitle{What Eclipse builds: \texttt{src} $\to$ \texttt{bin}}
  \lstterminal[basicstyle=\ttfamily\scriptsize]{codeoutput/Layout.txt}

  \vspace{0.1em}
  \small
  This is what \emph{Build Automatically} does on every save, done here by hand with the same compiler (\texttt{ecj}):
  one \texttt{.class} file per class, in the same package folders. Is the average right?
\end{frame}
% ~~~~~~~~~~~~~~~~~~~~~~~~~~~~~~~~~~~~~~ A frm




% ~~~~~~~~~~~~~~~~~~~~~~~~~~~~~~~~~~~~~~ V frm
\begin{frame}[t,fragile] \frametitle{The project files: \texttt{.project} and \texttt{.classpath}}
  \begin{columns}[T]
    \begin{column}{0.44\textwidth}
      \lstcode[basicstyle=\ttfamily\tiny]{codeoutput/Grades/.project}{XML}
    \end{column}
    \begin{column}{0.54\textwidth}
      \lstcode[basicstyle=\ttfamily\tiny]{codeoutput/Grades/.classpath}{XML}
    \end{column}
  \end{columns}
  \small \texttt{.project}: name and \emph{nature} (a Java project). \texttt{.classpath}: the \textbf{build path}: source folders,
  JRE, libraries, output folder. Edit it with \menu{Build Path > Configure Build Path\dots}, not by hand.
\end{frame}
% ~~~~~~~~~~~~~~~~~~~~~~~~~~~~~~~~~~~~~~ A frm




% ~~~~~~~~~~~~~~~~~~~~~~~~~~~~~~~~~~~~~~ V frm
\begin{frame}[t,fragile] \frametitle{The build path}
  \begin{center}\small
  \begin{tabular}{lll}
    \toprule
    \textbf{to use \dots}            & \textbf{in Configure Build Path}          & \textbf{becomes} \\
    \midrule
    a JAR file in the project      & \emph{Libraries} $\to$ \emph{Add JARs\dots}         & \texttt{kind="lib"} \\
    a JAR file elsewhere           & \emph{Add External JARs\dots}                       & \texttt{kind="lib"} (absolute path) \\
    JUnit 5                        & \emph{Add Library\dots} $\to$ JUnit                 & \texttt{kind="con"} \\
    another project of the workspace & \emph{Projects} $\to$ \emph{Add\dots}             & \texttt{kind="src"} \\
    another Java version           & \emph{JRE System Library} $\to$ \emph{Edit\dots}    & \texttt{kind="con"} \\
    \bottomrule
  \end{tabular}
  \end{center}
  \begin{itemize}\small
    \item Prefer JARs \textbf{inside} the project (e.g.\ a \texttt{lib/} folder): the project works on every computer.
          An external path like \texttt{C:\textbackslash Users\textbackslash ayse\textbackslash Downloads} does not.
    \item Real projects use Maven or Gradle to fetch libraries (junior level).
  \end{itemize}
\end{frame}
% ~~~~~~~~~~~~~~~~~~~~~~~~~~~~~~~~~~~~~~ A frm

% ~~~~~~~~~~~~~~~~~~~~~~~~~~~~~~~~~~~~~~ V
\hSection[Testing]{Unit Tests with JUnit 5}
% ~~~~~~~~~~~~~~~~~~~~~~~~~~~~~~~~~~~~~~ A




% ~~~~~~~~~~~~~~~~~~~~~~~~~~~~~~~~~~~~~~ V frm
\begin{frame}[t,fragile] \frametitle{A JUnit test class}
  \begin{columns}[T]
    \begin{column}{0.58\textwidth}
      \lstcode[basicstyle=\ttfamily\tiny]{codeoutput/p/GradeBookTest.java}{Java}
    \end{column}
    \begin{column}{0.38\textwidth}
      \begin{itemize}\small
        \item Right click the class under test: \menu{New > JUnit Test Case}, \emph{New JUnit Jupiter test}.
              Eclipse offers to add JUnit 5 to the build path.
        \item Keep tests in a separate source folder \texttt{test}, same package.
        \item \menu{Run As > JUnit Test}: the \textbf{JUnit view} shows a green or red bar.
      \end{itemize}
    \end{column}
  \end{columns}
\end{frame}
% ~~~~~~~~~~~~~~~~~~~~~~~~~~~~~~~~~~~~~~ A frm




% ~~~~~~~~~~~~~~~~~~~~~~~~~~~~~~~~~~~~~~ V frm
\begin{frame}[t,fragile] \frametitle{A red bar}
  \lstterminal[basicstyle=\ttfamily\scriptsize]{codeoutput/TestRed.txt}

  \vspace{0.1em}
  \small
  Same tests, run by the JUnit console launcher; the JUnit view in Eclipse shows the same tree. Expected 72.5, but was 72.0:
  where does the .5 go? Double-click the failure in the JUnit view to jump to the line.
\end{frame}
% ~~~~~~~~~~~~~~~~~~~~~~~~~~~~~~~~~~~~~~ A frm

% ~~~~~~~~~~~~~~~~~~~~~~~~~~~~~~~~~~~~~~ V
\hSection[Debugging]{Debugging}
% ~~~~~~~~~~~~~~~~~~~~~~~~~~~~~~~~~~~~~~ A




% ~~~~~~~~~~~~~~~~~~~~~~~~~~~~~~~~~~~~~~ V frm
\begin{frame}[t,fragile] \frametitle{Breakpoints that think}
  \begin{columns}[T]
    \begin{column}{0.5\textwidth}
      \begin{itemize}\small
        \item \textbf{Conditional breakpoint}: right click the blue dot, \emph{Breakpoint Properties\dots},
              \emph{Conditional}: \texttt{s.name().equals("Ali")}. It stops only then.
        \item \textbf{Hit count} 100: stops the 100th time, e.g.\ in a long loop.
        \item \textbf{Exception breakpoint}: \menu{Run > Add Java Exception Breakpoint\dots}, e.g.\
              \texttt{NullPointerException}: stops where it is thrown.
        \item \textbf{Watch}: select an expression, right click \emph{Watch}: it stays in the Expressions view.
      \end{itemize}
    \end{column}
    \begin{column}{0.46\textwidth}
      \begin{center}\footnotesize
      \begin{tabular}{@{}lll@{}}
        \toprule
        \textbf{action}     & \win                 & \mac \\
        \midrule
        toggle breakpoint   & \keys{Ctrl+\shift+B} & \keys{\cmd+\shift+B} \\
        debug last launch   & \keys{F11}           & \keys{\cmd+F11} \\
        inspect selection   & \keys{Ctrl+\shift+I} & \keys{\cmd+\shift+I} \\
        run to line         & \keys{Ctrl+R}        & \keys{\cmd+R} \\
        \bottomrule
      \end{tabular}
      \end{center}
      \small Stepping: \keys{F5} into, \keys{F6} over, \keys{F7} return, \keys{F8} resume.
    \end{column}
  \end{columns}
\end{frame}
% ~~~~~~~~~~~~~~~~~~~~~~~~~~~~~~~~~~~~~~ A frm




% ~~~~~~~~~~~~~~~~~~~~~~~~~~~~~~~~~~~~~~ V frm
\begin{frame}[t,fragile] \frametitle{Debugging the red test}
  \begin{columns}[T]
    \begin{column}{0.4\textwidth}
      \lstcode[basicstyle=\ttfamily\tiny, firstline=13, firstnumber=13]{codeoutput/p/GradeBook.java}{Java}
      \small In Eclipse: a breakpoint on line 18, \menu{Debug As > JUnit Test}, then look at the Variables view and
      type the expression into the Expressions view. Here: \texttt{jdb} with the commands of \texttt{debug.txt}.
    \end{column}
    \begin{column}{0.58\textwidth}
      \uncover<2->{\lstterminal[basicstyle=\ttfamily\tiny, firstline=8]{codeoutput/Debug.txt}}
    \end{column}
  \end{columns}
\end{frame}
% ~~~~~~~~~~~~~~~~~~~~~~~~~~~~~~~~~~~~~~ A frm




% ~~~~~~~~~~~~~~~~~~~~~~~~~~~~~~~~~~~~~~ V frm
\begin{frame}[t,fragile] \frametitle{The same steps with the command-line debugger}
  \begin{columns}[T]
    \begin{column}{0.48\textwidth}
      \begin{center}\footnotesize
      \begin{tabular}{@{}ll@{}}
        \toprule
        \textbf{Eclipse}                          & \textbf{\texttt{jdb}} \\
        \midrule
        breakpoint on line 18                    & \texttt{stop at \dots:18} \\
        \menu{Debug As}                          & \texttt{run} \\
        Variables view                           & \texttt{locals} \\
        Expressions view                         & \texttt{print \dots} \\
        \keys{F8} Resume                         & \texttt{cont} \\
        \keys{F6} Step Over                      & \texttt{next} \\
        \keys{F5} Step Into                      & \texttt{step} \\
        \bottomrule
      \end{tabular}
      \end{center}
    \end{column}
    \begin{column}{0.46\textwidth}
      \begin{itemize}\small
        \item The previous session is real: \texttt{jdb} talks to the program through the same JDWP as Eclipse.
        \item \texttt{slowly} types one command per second, as a person would.
        \item \texttt{sum / students.size()} is an \textbf{integer} division: \texttt{145 / 2 = 72}; the \texttt{.5} is lost before the result becomes \texttt{double}.
        \item Fix: \texttt{(double) sum / students.size()}.
      \end{itemize}
    \end{column}
  \end{columns}
\end{frame}
% ~~~~~~~~~~~~~~~~~~~~~~~~~~~~~~~~~~~~~~ A frm




% ~~~~~~~~~~~~~~~~~~~~~~~~~~~~~~~~~~~~~~ V frm
\begin{frame}[t,fragile] \frametitle{After the fix: a green bar}
  \lstterminal[basicstyle=\ttfamily\scriptsize]{codeoutput/TestGreen.txt}

  \vspace{0.1em}
  \small
  Re-run the tests after every fix: \keys{Ctrl+F11} / \keys{\cmd+\shift+F11} re-runs the last launch, including JUnit runs.
\end{frame}
% ~~~~~~~~~~~~~~~~~~~~~~~~~~~~~~~~~~~~~~ A frm

% ~~~~~~~~~~~~~~~~~~~~~~~~~~~~~~~~~~~~~~ V
\hSection[Refactoring]{Refactoring}
% ~~~~~~~~~~~~~~~~~~~~~~~~~~~~~~~~~~~~~~ A




% ~~~~~~~~~~~~~~~~~~~~~~~~~~~~~~~~~~~~~~ V frm
\begin{frame}[t,fragile] \frametitle{Before: it works, but \dots}
  \begin{columns}[T]
    \begin{column}{0.5\textwidth}
      \lstcode[basicstyle=\ttfamily\tiny]{codeoutput/Refactor/before/Report.java}{Java}
    \end{column}
    \begin{column}{0.46\textwidth}
      \begin{itemize}\small
        \item Names say nothing: \texttt{g}, \texttt{s}, \texttt{c}, \texttt{x}.
        \item One long \texttt{main} doing two jobs.
        \item A magic number: \texttt{50}.
      \end{itemize}
      \vspace{0.3em}
      \begin{center}\scriptsize
      \begin{tabular}{@{}lll@{}}
        \toprule
        \textbf{refactoring}  & \win                & \mac \\
        \midrule
        Rename                & \keys{Alt+\shift+R} & \keys{Option+\cmd+R} \\
        Extract Method        & \keys{Alt+\shift+M} & \keys{Option+\cmd+M} \\
        Extract Constant      & \menu{Refactor}     & \menu{Refactor} \\
        \bottomrule
      \end{tabular}
      \end{center}
    \end{column}
  \end{columns}
\end{frame}
% ~~~~~~~~~~~~~~~~~~~~~~~~~~~~~~~~~~~~~~ A frm




% ~~~~~~~~~~~~~~~~~~~~~~~~~~~~~~~~~~~~~~ V frm
\begin{frame}[t,fragile] \frametitle{After: three refactorings}
  \begin{columns}[T]
    \begin{column}{0.5\textwidth}
      \lstcode[basicstyle=\ttfamily\tiny]{codeoutput/Refactor/after/Report.java}{Java}
    \end{column}
    \begin{column}{0.46\textwidth}
      \begin{enumerate}\small
        \item \textbf{Rename} \texttt{g}, \texttt{s}, \texttt{c}, \texttt{x}: every use changes at once.
        \item Select a loop, \textbf{Extract Method}: Eclipse finds the parameters and the return value.
        \item Select \texttt{50}, \textbf{Extract Constant}: \texttt{PASS\_GRADE}.
      \end{enumerate}
      \uncover<2->{\lstterminal[basicstyle=\ttfamily\scriptsize]{codeoutput/Refactor.txt}
      \small Same output before and after: that is what \emph{refactoring} means.}
    \end{column}
  \end{columns}
\end{frame}
% ~~~~~~~~~~~~~~~~~~~~~~~~~~~~~~~~~~~~~~ A frm




% ~~~~~~~~~~~~~~~~~~~~~~~~~~~~~~~~~~~~~~ V frm
\begin{frame}[t,fragile] \frametitle{Let Eclipse write the routine code}
  \begin{center}\small
  \begin{tabular}{ll}
    \toprule
    \textbf{menu \menu{Source > \dots}}                  & \textbf{writes} \\
    \midrule
    Generate Constructor using Fields\dots             & a constructor that sets the chosen fields \\
    Generate Getters and Setters\dots                  & \texttt{getName()}, \texttt{setName(\dots)} \\
    Generate \texttt{toString()}\dots                  & a \texttt{toString} listing the fields \\
    Generate \texttt{hashCode()} and \texttt{equals()}\dots & both, consistent with each other \\
    Override/Implement Methods\dots                    & stubs for interface or superclass methods \\
    \bottomrule
  \end{tabular}
  \end{center}
  \begin{itemize}\small
    \item Right click in the editor, or \keys{Alt+\shift+S} (\win) / \keys{Option+\cmd+S} (\mac).
    \item For pure data classes, a Java \texttt{record} (like \texttt{Student}) gives all of these without any code.
  \end{itemize}
\end{frame}
% ~~~~~~~~~~~~~~~~~~~~~~~~~~~~~~~~~~~~~~ A frm

% ~~~~~~~~~~~~~~~~~~~~~~~~~~~~~~~~~~~~~~ V
\hSection[Git]{Git in Eclipse}
% ~~~~~~~~~~~~~~~~~~~~~~~~~~~~~~~~~~~~~~ A




% ~~~~~~~~~~~~~~~~~~~~~~~~~~~~~~~~~~~~~~ V frm
\begin{frame}[t,fragile] \frametitle{The Team menu is Git}
  \begin{columns}[T]
    \begin{column}{0.52\textwidth}
      \begin{center}\small
      \begin{tabular}{ll}
        \toprule
        \textbf{Eclipse (EGit)}                       & \textbf{git} \\
        \midrule
        \menu{Team > Share Project\dots}             & \texttt{git init} \\
        Git Staging view: drag to \emph{Staged}      & \texttt{git add} \\
        \emph{Commit} button                         & \texttt{git commit} \\
        \menu{Team > Pull}, \menu{Team > Push\dots}  & \texttt{git pull}, \texttt{git push} \\
        \menu{Team > Show in History}                & \texttt{git log} \\
        \menu{Team > Switch To}                      & \texttt{git switch} \\
        \menu{File > Import > Git > Projects from Git} & \texttt{git clone} \\
        \bottomrule
      \end{tabular}
      \end{center}
    \end{column}
    \begin{column}{0.44\textwidth}
      \begin{itemize}\small
        \item EGit is part of the Java package. Files show their state with small icons
              (\texttt{>} changed, \texttt{?} untracked).
        \item Commit \texttt{.project}, \texttt{.classpath}, \texttt{.settings/}: they \emph{are} the project.
        \item Never commit \texttt{bin/} (built) or the workspace \texttt{.metadata/} (personal).
      \end{itemize}
    \end{column}
  \end{columns}
\end{frame}
% ~~~~~~~~~~~~~~~~~~~~~~~~~~~~~~~~~~~~~~ A frm




% ~~~~~~~~~~~~~~~~~~~~~~~~~~~~~~~~~~~~~~ V frm
\begin{frame}[t,fragile] \frametitle{An Eclipse project in Git}
  \begin{columns}[T]
    \begin{column}{0.49\textwidth}
      \lstterminal[basicstyle=\ttfamily\tiny, lastline=17]{codeoutput/Git.txt}
    \end{column}
    \begin{column}{0.49\textwidth}
      \lstterminal[basicstyle=\ttfamily\tiny, firstline=18]{codeoutput/Git.txt}
    \end{column}
  \end{columns}

  \vspace{0.1em}
  \small
  One line of \texttt{.gitignore} keeps the compiled classes out; the project files stay in, so a classmate can
  \emph{Import} the project and build it at once.
\end{frame}
% ~~~~~~~~~~~~~~~~~~~~~~~~~~~~~~~~~~~~~~ A frm

% ~~~~~~~~~~~~~~~~~~~~~~~~~~~~~~~~~~~~~~ V
\hSection[Export]{Running Outside Eclipse}
% ~~~~~~~~~~~~~~~~~~~~~~~~~~~~~~~~~~~~~~ A




% ~~~~~~~~~~~~~~~~~~~~~~~~~~~~~~~~~~~~~~ V frm
\begin{frame}[t,fragile] \frametitle{Export a runnable JAR}
  \begin{columns}[T]
    \begin{column}{0.46\textwidth}
      \begin{enumerate}\small
        \item \menu{File > Export > Java > Runnable JAR file}.
        \item \emph{Launch configuration}: \texttt{Main - Grades} (made by the first run).
        \item \emph{Export destination}: \texttt{grades.jar}. \textbf{Finish}.
        \item Run it anywhere with Java:\\ \hCode{java -jar grades.jar}
      \end{enumerate}
      \small The JAR's manifest names the main class, so \texttt{java -jar} knows where to start.
    \end{column}
    \begin{column}{0.5\textwidth}
      \uncover<2->{\lstterminal[basicstyle=\ttfamily\scriptsize]{codeoutput/Jar.txt}
      \small The same JAR made with the JDK's \texttt{jar} tool (Command Prompt on \win, Terminal on \mac).}
    \end{column}
  \end{columns}
\end{frame}
% ~~~~~~~~~~~~~~~~~~~~~~~~~~~~~~~~~~~~~~ A frm

% ~~~~~~~~~~~~~~~~~~~~~~~~~~~~~~~~~~~~~~ V
\hSection[Habits]{Good Habits}
% ~~~~~~~~~~~~~~~~~~~~~~~~~~~~~~~~~~~~~~ A




% ~~~~~~~~~~~~~~~~~~~~~~~~~~~~~~~~~~~~~~ V frm
\begin{frame}[t,fragile] \frametitle{Habits that save hours}
  \begin{columns}[T]
    \begin{column}{0.48\textwidth}
      \begin{block}{Every save}
        \small \textbf{Save Actions}: \menu{Window > Preferences > Java > Editor > Save Actions} (\mac: \menu{Eclipse > Settings\dots}):
        format edited lines, organize imports.
      \end{block}
      \begin{block}{Every change}
        \small Run the tests (\menu{Run As > JUnit Test}); keep the Problems view at zero errors and few warnings.
      \end{block}
    \end{column}
    \begin{column}{0.48\textwidth}
      \begin{block}{Every bug}
        \small Write a failing test first, debug, fix, see it turn green. The test keeps the bug from coming back.
      \end{block}
      \begin{block}{Every day}
        \small Commit small steps with clear messages; push to the course repository.
      \end{block}
    \end{column}
  \end{columns}
\end{frame}
% ~~~~~~~~~~~~~~~~~~~~~~~~~~~~~~~~~~~~~~ A frm

% ~~~~~~~~~~~~~~~~~~~~~~~~~~~~~~~~~~~~~~ V
\hSection[Summary]{Summary}
% ~~~~~~~~~~~~~~~~~~~~~~~~~~~~~~~~~~~~~~ A




% ~~~~~~~~~~~~~~~~~~~~~~~~~~~~~~~~~~~~~~ V frm
\begin{frame}[t,fragile] \frametitle{Cheat sheet}
  \begin{center}\footnotesize
  \begin{tabular}{@{}lll@{}}
    \toprule
    \textbf{task}           & \textbf{menu}                                   & \textbf{\win{} / \mac} \\
    \midrule
    new package / class     & \menu{File > New > Package}, \menu{Class}       & \\
    build path              & \menu{Build Path > Configure Build Path\dots}   & \\
    new test                & \menu{File > New > JUnit Test Case}             & \\
    run tests               & \menu{Run As > JUnit Test}                      & \keys{Alt+\shift+X}, \keys{T} / \keys{Option+\cmd+X}, \keys{T} \\
    toggle breakpoint       & \menu{Run > Toggle Breakpoint}                  & \keys{Ctrl+\shift+B} / \keys{\cmd+\shift+B} \\
    rename                  & \menu{Refactor > Rename\dots}                   & \keys{Alt+\shift+R} / \keys{Option+\cmd+R} \\
    extract method          & \menu{Refactor > Extract Method\dots}           & \keys{Alt+\shift+M} / \keys{Option+\cmd+M} \\
    generate code           & \menu{Source}                                   & \keys{Alt+\shift+S} / \keys{Option+\cmd+S} \\
    runnable JAR            & \menu{File > Export > Runnable JAR file}        & \\
    \bottomrule
  \end{tabular}
  \end{center}
\end{frame}
% ~~~~~~~~~~~~~~~~~~~~~~~~~~~~~~~~~~~~~~ A frm




% ~~~~~~~~~~~~~~~~~~~~~~~~~~~~~~~~~~~~~~ V frm
\begin{frame}[t,fragile] \frametitle{Exercises}
  \begin{enumerate}\small
    \item Create the project \texttt{Grades} with the package \texttt{tr.edu.yeditepe.grades} and the three classes of this deck.
    \item Add the JUnit test class, see the red bar, debug with a breakpoint on line 18 and fix the bug.
    \item Add a method \texttt{countPassed()} to \texttt{GradeBook}: first the test, then the code.
    \item Refactor \texttt{Report} yourself with Rename, Extract Method and Extract Constant; compare with this deck.
    \item Put the project into Git with EGit, push it to the course server, and import it on another computer.
  \end{enumerate}
\end{frame}
% ~~~~~~~~~~~~~~~~~~~~~~~~~~~~~~~~~~~~~~ A frm




% ~~~~~~~~~~~~~~~~~~~~~~~~~~~~~~~~~~~~~~ V frm
\begin{frame}[t,fragile]
	\begin{center}
		\vspace{4cm}
		\hRemark{\Huge Questions?}
	\end{center}
\end{frame}
% ~~~~~~~~~~~~~~~~~~~~~~~~~~~~~~~~~~~~~~ A frm



\end{document}
